
\chapter{Aplicaci�n de Data Discovery a datos de instituciones del Estado}

\begin{wrapfigure}{R}{0pt}
	\centering
	\includegraphics[width=0.25\textwidth]{figuras/makeDecision}
	\caption{}
	\label{fig:makeDecision}
\end{wrapfigure}

En Paraguay, e inclusive mundialmente, la mayor�a de las empresas no logran comprender 100% los datos que generan. La consecuencia de no comprender los datos puede conllevar a tomar una decisi�n equivocada y esa decisi�n puede terminar en algo irreversible y de mucho impacto negativo para la organizaci�n a lo largo del tiempo. La informaci�n es considerada como uno de los recursos m�s importantes en una empresa, porque en base a esto se obtiene el conocimiento.

Aplicando la t�cnica de data discovery podremos detectar irregularidades, predecir acontecimientos futuros y preverlos a tiempo. En este trabajo analizaremos los datos de la ANDE y de la DGEEC, realizando cruzamiento entre ambos, con el objetivo de obtener informaci�n de inter�s para la entidad. 

\subsection{Datos de la ANDE y de la DGEEC}

Se cuenta con datos de consumo de energ�a el�ctrica, importe facturado, grupo de consumo (residencial, industrial, exportaci�n, etc), por a�o (2000-2014), departamento y distrito. Estos datos fueron solicitados formalmente a la entidad a trav�s de la Facultad de Ciencias y tecnolog�a de la Universidad Cat�lica, la cual tuvimos una respuesta favorable para proceder.

\subsection{Dashboard de control / monitoramiento}
	
En esta secci�n se mostrar�n 4 ejemplos de paneles hechos con los datos de ambas instituciones (DGEEC y ANDE).  Una de las t�cnicas de comparaci�n a ser utilizadas en los gr�ficos es ,  ?Medici�n de tasas de crecimientos?, la cual se calcula el porcentaje de crecimiento que hubo por cada a�o. Por ejemplo, si al cerrar el a�o 2014, la cantidad de clientes lleg� a 1.000.000 y en el a�o 2015 aument� 100.000, esto quiere decir que en el a�o 2015, la tasa de crecimiento de los clientes fue del 10%, es decir, hubo un crecimiento positivo y la cantidad de clientes ha aumentado el respecto al a�o anterior. Suponiendo que el a�o 2016 la ANDE cierra con un total de 1.000.000 clientes,, su crecimiento ser� de 10% negativo.
La f�rmula empleada es la que sigue, donde ?n? es el a�o actual y ?n-1? el a�o anterior, PIB es una variable que indica, en el caso  de nuestra comparaci�n, la cantidad de clientes que posee la ANDE .
	
\section{Data marts}
Escribir sobre los objetivos
Como puede ser visto en la figura \ref{fig:BI}.
\section{Data warehouse}